% =============================================================================
% Module 9: Solutions to Exercises
% =============================================================================
% This file is conditionally included based on \ifsolutions
% =============================================================================

\section{Solutions to Exercises}
\label{sec:galaxy_lensing_solutions}

\noindent
Full worked solutions are also available in the Mathematica script
\solnlink{09_Galaxy_Lensing_Applications/problems\_09.wl}{problems\_09.wl},
which verifies each answer symbolically and numerically.

% ----- Exercise 9.1 -----
\subsection*{Exercise~\ref{ex:sie_mass}: Enclosed Mass for an SIE Lens}

\textbf{(a)} For a flat $\Lambda$CDM cosmology with $\Hubble = 70$~km/s/Mpc
and $\Omegam = 0.3$, the angular diameter distances are computed via:
\begin{equation*}
    D(z_1, z_2) = \frac{\speedoflight}{(1+z_2)\,\Hubble}
    \int_{z_1}^{z_2} \frac{dz'}{E(z')}
\end{equation*}
where $E(z) = \sqrt{\Omegam(1+z)^3 + \OmegaL}$.  Numerical
evaluation gives:
\begin{align*}
    \Dd &\approx 770~\mathrm{Mpc} \,, \\
    \Ds &\approx 1640~\mathrm{Mpc} \,, \\
    \Dds &\approx 1180~\mathrm{Mpc} \,.
\end{align*}
(Exact values depend on the cosmological parameters and the
integration method; see \texttt{problems\_09.wl} for precise values.)

\textbf{(b)} The critical surface mass density is:
\begin{equation*}
    \Sigmacr = \frac{\speedoflight^2}{4\pi \Gnewton}
    \,\frac{\Ds}{\Dd\,\Dds}
    \approx 3.2 \times 10^9\;\Msun\,\mathrm{kpc}^{-2} \,.
\end{equation*}

\textbf{(c)} The enclosed projected mass is:
\begin{equation*}
    M_{\mathrm{E}} = \pi\,\thetaE^2\,\Dd^2\,\Sigmacr
    = \frac{\speedoflight^2}{4\Gnewton}\,\thetaE^2\,
    \frac{\Dd\,\Ds}{\Dds} \,.
\end{equation*}
Converting $\thetaE = 1.5'' = 7.27 \times 10^{-6}$~rad:
\begin{equation*}
    M_{\mathrm{E}} \approx 4.2 \times 10^{11}\;\Msun \,.
\end{equation*}
This is larger than the typical stellar mass of a massive elliptical
($\sim 10^{11}\;\Msun$), indicating significant dark matter within
the Einstein radius. \checkmark

\textbf{(d)} From $\thetaE = 4\pi\sigma_v^2\Dds/(\speedoflight^2\Ds)$:
\begin{equation*}
    \sigma_v = \speedoflight\,\sqrt{\frac{\thetaE\,\Ds}{4\pi\,\Dds}}
    \approx 270~\mathrm{km/s} \,.
\end{equation*}
This is typical of a massive early-type galaxy. \checkmark

% ----- Exercise 9.2 -----
\subsection*{Exercise~\ref{ex:msd_derivation}: Mass-Sheet Degeneracy}

\textbf{(a)} The original deflection satisfies
$\nabla \cdot \vecalpha = 2\kappab$.  Under the MST:
\begin{equation*}
    \kappab_\lambda = \lambda\kappab + (1-\lambda)
    \implies
    \vecalpha_\lambda = \lambda\vecalpha + (1-\lambda)\vectheta \,.
\end{equation*}
The transformed lens equation:
\begin{align*}
    \vecbeta_\lambda
    &= \vectheta - \vecalpha_\lambda(\vectheta) \\
    &= \vectheta - \lambda\vecalpha(\vectheta)
       - (1-\lambda)\vectheta \\
    &= \lambda[\vectheta - \vecalpha(\vectheta)] \\
    &= \lambda\vecbeta \,.
\end{align*}
So the same image positions $\vectheta$ satisfy the transformed lens
equation for source position $\lambda\vecbeta$. \checkmark

\textbf{(b)} The Jacobian is
$\Amat = \mathbb{I} - \partial\vecalpha/\partial\vectheta$.
Under the MST:
\begin{align*}
    \Amat_\lambda
    &= \mathbb{I}
       - \frac{\partial\vecalpha_\lambda}{\partial\vectheta} \\
    &= \mathbb{I} - \lambda\frac{\partial\vecalpha}{\partial\vectheta}
       - (1-\lambda)\mathbb{I} \\
    &= \lambda\left(\mathbb{I}
       - \frac{\partial\vecalpha}{\partial\vectheta}\right)
    = \lambda\,\Amat \,.
\end{align*}
Therefore $\det\Amat_\lambda = \lambda^2\det\Amat$, and
$\mu_\lambda = 1/\det\Amat_\lambda = \mu/\lambda^2$. \checkmark

\textbf{(c)} Since $\mu_{\lambda,i} = \mu_i/\lambda^2$ for every
image $i$, the ratio:
\begin{equation*}
    \frac{\mu_{\lambda,i}}{\mu_{\lambda,j}}
    = \frac{\mu_i/\lambda^2}{\mu_j/\lambda^2}
    = \frac{\mu_i}{\mu_j} \,.
\end{equation*}
The $\lambda^2$ factors cancel.  Flux ratios are invariant. \checkmark

\textbf{(d)} The Fermat potential is
$\phi(\vectheta, \vecbeta) = \frac{1}{2}|\vectheta - \vecbeta|^2
- \psiL(\vectheta)$.  Under the MST, the lensing potential
transforms as $\psiL_\lambda = \lambda\psiL +
\frac{1}{2}(1-\lambda)|\vectheta|^2$.  The time delay between
images $A$ and $B$ is:
\begin{align*}
    \Delta t_\lambda
    &= \frac{(1+z_d)\Dd\Ds}{\speedoflight\,\Dds}
       [\phi_\lambda(\vectheta_A, \vecbeta_\lambda)
        - \phi_\lambda(\vectheta_B, \vecbeta_\lambda)] \\
    &= \lambda\,\Delta t \,.
\end{align*}
The derivation proceeds by substituting $\vecbeta_\lambda =
\lambda\vecbeta$ and $\psiL_\lambda = \lambda\psiL +
\frac{1}{2}(1-\lambda)|\vectheta|^2$ and simplifying.  The
uniform $|\vectheta|^2$ term cancels between the two images
because they are at the same image positions, leaving a factor
of $\lambda$ multiplying the original Fermat potential
difference. \checkmark

% ----- Exercise 9.3 -----
\subsection*{Exercise~\ref{ex:H0_time_delay}: $\Hubble$ from a Time-Delay Lens}

\textbf{(a)} The time-delay equation is:
\begin{equation*}
    \Delta t = \frac{(1+z_d)}{\speedoflight}
    \,\frac{\Dd\,\Ds}{\Dds}\,\Delta\tau \,.
\end{equation*}
Solving for $\Hubble$: since $\Dd\Ds/\Dds \propto 1/\Hubble$,
we can write
$\Dd\Ds/\Dds = d(z_d, z_s, \Omegam)/\Hubble$ where $d$ is a
dimensionless function of redshifts and cosmological parameters.
Then:
\begin{equation*}
    \Hubble = \frac{(1+z_d)\,\Delta\tau}{\speedoflight\,\Delta t}
    \times d(z_d, z_s, \Omegam) \,.
\end{equation*}

\textbf{(b)} For $z_d = 0.6$, $z_s = 1.5$, $\Omegam = 0.3$
(flat $\Lambda$CDM), the dimensionless distance ratio
$d = \Hubble\Dd\Ds/\Dds$ is computed from the cosmological
integrals.  Converting $\Delta\tau = 0.64~\mathrm{arcsec}^2
= 0.64 \times (4.848 \times 10^{-6})^2~\mathrm{rad}^2
= 1.50 \times 10^{-11}~\mathrm{rad}^2$ and
$\Delta t = 80~\mathrm{days} = 6.91 \times 10^6~\mathrm{s}$:

Numerical evaluation gives $\Hubble \approx 70$~km/s/Mpc
(the exact value depends on the precise cosmological distances;
see \texttt{problems\_09.wl} for the detailed calculation). \checkmark

\textbf{(c)} If the true lens has an unmodeled mass sheet with
$\kappab_s = 0.05$, then $\lambda = 1 - \kappab_s = 0.95$.  The
inferred $\Hubble$ is biased:
\begin{equation*}
    \Hubble^{\mathrm{inferred}} = \lambda\,\Hubble^{\mathrm{true}}
    = 0.95\,\Hubble^{\mathrm{true}} \,.
\end{equation*}
The inferred $\Hubble$ would be 5\% too low.  For a true value of
$\sim$73~km/s/Mpc, the inferred value would be $\sim$69~km/s/Mpc
--- close to the Planck value.  This illustrates why the MSD is
central to the Hubble tension discussion. \checkmark

% ----- Exercise 9.4 -----
\subsection*{Exercise~\ref{ex:power_law_profile}: Power-Law Profile}

\textbf{(a)} For $\rho(r) = \rho_0(r/r_0)^{-\gamma'}$, the projected
surface mass density is:
\begin{equation*}
    \Sigma(R) = \int_{-\infty}^{\infty} \rho(\sqrt{R^2 + z^2})\,dz
    = \rho_0\,r_0^{\gamma'}\,\frac{\sqrt{\pi}\,\Gamma[(\gamma'-1)/2]}
    {\Gamma[\gamma'/2]}\,R^{1-\gamma'}
\end{equation*}
where $R = \Dd\theta$ is the projected physical radius.  The
convergence is $\kappab = \Sigma/\Sigmacr$, so:
\begin{equation*}
    \kappab(\theta) \propto \theta^{1-\gamma'} \,.
\end{equation*}
More precisely, $\kappab(\theta) = \frac{3-\gamma'}{2}
(\thetaE/\theta)^{\gamma'-1}$, where $\thetaE$ is defined as the
radius where the mean convergence equals unity. \checkmark

\textbf{(b)} For $\gamma' = 2$:
$\kappab \propto \theta^{1-2} = \theta^{-1} = 1/\theta$, which
is the SIS convergence $\kappab = \thetaE/(2\theta)$. \checkmark

\textbf{(c)} The enclosed mass within radius $R$ scales as
$M(\le R) \propto R^{3-\gamma'}$.  The Einstein mass is
$M_{\mathrm{E}} = M(\le R_{\mathrm{E}})$, fixed by lensing.
The velocity dispersion scales as:
\begin{equation*}
    \sigma_v^2 \propto \frac{M(\le R_{\mathrm{ap}})}{R_{\mathrm{ap}}}
    \propto R_{\mathrm{ap}}^{2-\gamma'} \,.
\end{equation*}
Since $R_{\mathrm{ap}} \neq R_{\mathrm{E}}$ in general, the ratio
$\sigma_v^2/M_{\mathrm{E}}$ depends on $\gamma'$, breaking the
degeneracy between mass normalization and profile slope. \checkmark

\textbf{(d)} For a given $\thetaE$ and $\sigma_v$, the self-consistent
$\gamma'$ can be found from the relation linking both.  For
$\thetaE = 1.2''$ and $\sigma_v = 230$~km/s, numerical evaluation
(see \texttt{problems\_09.wl}) gives $\gamma' \approx 2.0$--$2.1$,
fully consistent with the bulge-halo conspiracy value
$\gamma' \approx 2.0 \pm 0.2$. \checkmark

% ----- Exercise 9.5 -----
\subsection*{Exercise~\ref{ex:compare_lenses}: Comparing Real Lens Systems}

\textbf{(a)} Angular diameter distances (flat $\Lambda$CDM,
$\Hubble = 70$~km/s/Mpc, $\Omegam = 0.3$):

\begin{center}
\begin{tabular}{lccc}
\toprule
\textbf{System} & $\Dd$ (Mpc) & $\Ds$ (Mpc) & $\Dds$ (Mpc) \\
\midrule
Q0957+561  & $\sim$1040 & $\sim$1740 & $\sim$1155 \\
Q2237+0305 & $\sim$163  & $\sim$1747 & $\sim$1683 \\
B1608+656  & $\sim$1410 & $\sim$1739 & $\sim$777  \\
\bottomrule
\end{tabular}
\end{center}

\textbf{(b)} Enclosed projected masses
($M_{\mathrm{E}} = \speedoflight^2\thetaE^2\Dd\Ds/(4\Gnewton\Dds)$):

\begin{center}
\begin{tabular}{lc}
\toprule
\textbf{System} & $M_{\mathrm{E}}$ ($10^{11}\,\Msun$) \\
\midrule
Q0957+561  & $\sim 1.9$ \\
Q2237+0305 & $\sim 0.2$ \\
B1608+656  & $\sim 3.9$ \\
\bottomrule
\end{tabular}
\end{center}

\textbf{(c)} Physical Einstein radii ($R_{\mathrm{E}} = \Dd\,\thetaE$):

\begin{center}
\begin{tabular}{lc}
\toprule
\textbf{System} & $R_{\mathrm{E}}$ (kpc) \\
\midrule
Q0957+561  & $\sim 5.0$ \\
Q2237+0305 & $\sim 0.7$ \\
B1608+656  & $\sim 6.8$ \\
\bottomrule
\end{tabular}
\end{center}

\textbf{(d)} B1608+656 has both the largest enclosed mass ($\sim 3.9 \times
10^{11}\,\Msun$) and the largest physical Einstein radius
($\sim 6.8$~kpc), despite having a similar angular Einstein
radius to Q0957+561.  This is because B1608+656 is at higher
redshift ($z_d = 0.63$), so the same angular size corresponds to
a larger physical size and a larger mass.  Q2237+0305 has the
smallest mass and physical radius because its lens is so nearby
($z_d = 0.04$) --- the angular diameter distance $\Dd$ is small,
so the physical Einstein radius is only $\sim 0.7$~kpc even though
$\thetaE \approx 0.9''$. \checkmark
